\section*{الفصل \ref{ch:b3:complex-spectra-magnetism} --- الأطياف الإلكترونية للمعقدات ومغناطيسيتها}

\begin{solution}{exo:b3:complex-spectra-magnetism:1}
D: \babelsublr{E$_g$ + T$_{2g}$} ($2 + 3 = 5$)؛ F: \babelsublr{A$_{2g}$ + T$_{1g}$ + T$_{2g}$} ($1 + 3 + 3 = 7$)؛ G: \babelsublr{A$_{1g}$ + E$_g$ + T$_{1g}$ + T$_{2g}$}
($1 + 2 + 3 + 3 = 9$).
\end{solution}

\begin{solution}{exo:b3:complex-spectra-magnetism:2}
\babelsublr{$d^1$ $^2$D $\to$ $^2$T$_{2g}$}؛ \babelsublr{$d^2$ $^3$F $\to$ $^3$T$_{1g}$}؛ \babelsublr{$d^3$ $^4$F $\to$ $^4$A$_{2g}$}؛ \babelsublr{$d^4$ $^5$D $\to$ $^5$E$_g$}؛ \babelsublr{$d^5$ $^6$S $\to$ $^6$A$_{1g}$}؛ \babelsublr{$d^6$ $^5$D
$\to$ $^5$T$_{2g}$}؛ \babelsublr{$d^7$ $^4$F $\to$ $^4$T$_{1g}$}؛ \babelsublr{$d^8$ $^3$F $\to$ $^3$A$_{2g}$}؛ \babelsublr{$d^9$ $^2$D $\to$ $^2$E$_g$}.
\end{solution}

\begin{solution}{exo:b3:complex-spectra-magnetism:3}
عالي السبين: \ce{Mn^{2+}} ($n = 5$) $5.92\,\mu_B$، \ce{Fe^{2+}} (4) $4.90\,\mu_B$، \ce{Co^{2+}} (3) $3.87\,\mu_B$؛ ومنخفض السبين: \ce{Fe^{2+}} 
و\ce{Co^{3+}} ($t_{2g}^6$): 0، ديامغناطيسيان.
\end{solution}

\begin{solution}{exo:b3:complex-spectra-magnetism:4}
$\ce{[Mn(H2O)6]^{2+}}$ (ممنوع سبينيًا ووفق لابورت، شبه عديم اللون) $<$ $\ce{[Co(H2O)6]^{2+}}$
(ممنوع وفق لابورت، \omterm{def:b3:electronic-spectroscopy:vibronic}{اهتزازي إلكتروني}) $<$ $\ce{[CoCl4]^{2-}}$ (لا مركز تناظر) $<$ \ce{MnO4-} (نقل شحنة
مسموح).
\end{solution}

\begin{solution}{exo:b3:complex-spectra-magnetism:5}
8500: \babelsublr{$^3$A$_{2g}\to{}^3$T$_{2g}$}، ومنه $\Delta_{\mathrm o} = \qty{8500}{cm^{-1}}$؛ 14100: \babelsublr{$^3$T$_{1g}$(F)}؛ 24900: \babelsublr{$^3$T$_{1g}$(P)}.
$15B = 39\,000 - 3 \times 8500$، $B = \qty{900}{cm^{-1}}$، $\beta = 900/1056 = 0.85$.
\end{solution}

\begin{solution}{exo:b3:complex-spectra-magnetism:6}
$\Delta_{\mathrm o} = \qty{17400}{cm^{-1}}$؛ وبحل $7.5B + 1.5\Delta_{\mathrm o} - \frac12\sqrt{225B^2 - 18B\Delta_{\mathrm o} + \Delta_{\mathrm o}^2} = 24\,600$: $B = \qty{729}{cm^{-1}}$،
$\beta = 0.79$. $\nu_3 = \qty{38500}{cm^{-1}}$ (\qty{260}{nm})، في فوق البنفسجي.
\end{solution}

\begin{solution}{exo:b3:complex-spectra-magnetism:7}
إلكترونات $d$ أكثر انتشارًا على أيونات الأكسيد في الحجرين منها على
جزيئات الماء: فالروابط Cr--O فيهما أكثر تساهمية.
\end{solution}

\begin{solution}{exo:b3:complex-spectra-magnetism:8}
قوي ($e_g$ مملوءة على نحو غير متساوٍ): $d^9$، و$d^4$ عالي السبين، و$d^7$ منخفض السبين. منعدم: $d^3$ ($t_{2g}^3$، حد A) 
و$d^6$ منخفض السبين ($t_{2g}^6$). ضعيف: $d^6$ عالي السبين ($t_{2g}^4e_g^2$، حد T).
\end{solution}

\begin{solution}{exo:b3:complex-spectra-magnetism:9}
$d^7$ ثماني الأوجه عالي السبين: الحد الأساسي \babelsublr{$^4$T$_{1g}$}، حد T ذو \omterm{def:b3:complex-spectra-magnetism:moment}{إسهام مداري}،
والعزم أعلى كثيرًا من قيمة السبين وحده $3.87\,\mu_B$. و$d^7$ رباعي الأوجه ($\equiv$ $d^3$ ثماني الأوجه): الحد
الأساسي \babelsublr{$^4$A$_2$}، والعزم المداري مُخمَد، قريب من قيمة السبين وحده (أعلى قليلًا، بسبب
الامتزاج مع حدود T المثارة).
\end{solution}

\begin{solution}{exo:b3:complex-spectra-magnetism:10}
$\mu_{\mathrm{eff}} = 2.828\sqrt{1.87} = 3.87\,\mu_B$: ثلاثة إلكترونات منفردة ($S = \frac32$).
\end{solution}

\begin{solution}{exo:b3:complex-spectra-magnetism:11}
$\Delta\chi_v = 3 \times 25.0/\num{400e6} = \num{1.88e-7}$؛ $\chi_{\mathrm m} = \num{1.88e-7}/\qty{10.0}{mol.m^{-3}} = \qty{1.88e-8}{m^3.mol^{-1}}$، أي
\qty{1.49e-3}{cm^3.mol^{-1}}؛ $\chi_{\mathrm m}T = 0.445$، $\mu_{\mathrm{eff}} = 1.89\,\mu_B$: إلكترون منفرد واحد.
\end{solution}

\begin{solution}{exo:b3:complex-spectra-magnetism:12}
$T_{1/2} = 15\,000/75 = \qty{200}{K}$. وعند \qty{250}{K}: $\Delta G = 15\,000 - 250 \times 75 = \qty{-3750}{J/mol}$، $K = \eu^{1.80} = 6.07$،
$x_{\mathrm{HS}} = 0.86$؛ $\chi_{\mathrm m}T = 0.86 \times 3.00 = \qty{2.6}{cm^3.K.mol^{-1}}$.
\end{solution}

\begin{solution}{pb:b3:complex-spectra-magnetism:1}
\textbf{1.} \babelsublr{[Ar]$3d^3$}.
\textbf{2.} \babelsublr{$^4$F}.
\textbf{3.} \babelsublr{$^4$A$_{2g}$ + $^4$T$_{2g}$ + $^4$T$_{1g}$}.
\textbf{4.} \babelsublr{$^4$A$_{2g}$}، حد التشكيل $t_{2g}^3$، والإلكترونات الثلاثة كلها في المدارات الدنيا.
\textbf{5.} \babelsublr{$^4$T$_{1g}$}(P).
\textbf{6.} \babelsublr{$^4$A$_{2g}\to{}^4$T$_{2g}$, $\to{}^4$T$_{1g}$(F), $\to{}^4$T$_{1g}$(P)}.
\textbf{7.} 561 و\qty{414}{nm}.
\textbf{8.} 597 و\qty{434}{nm}.
\textbf{9.} يمتص الياقوت الأصفر المخضر والبنفسجي وينفذ منه الأحمر (مع قليل من
الأزرق)؛ وحزم الزمرد، المنزاحة نحو الأحمر، تمتص الأحمر البرتقالي أيضًا، فينفذ
الأخضر.
\textbf{10.} الياقوت \qty{17820}{cm^{-1}}، والزمرد \qty{16740}{cm^{-1}}.
\textbf{11.} للحدين \babelsublr{$^4$T$_{2g}$} ($t_{2g}^2e_g$) و\babelsublr{$^4$A$_{2g}$} ($t_{2g}^3$) طاقة التنافر الإلكتروني نفسها؛ 
والفرق بينهما هو طاقة الترقية الأحادية الإلكترون $\Delta_{\mathrm o}$.
\textbf{12.} $B = (14\,305 - 547)/15 = \qty{917}{cm^{-1}}$.
\textbf{13.} \qty{614}{cm^{-1}}.
\textbf{14.} \qty{618}{cm^{-1}}.
\textbf{15.} 0.67 لكليهما.
\textbf{16.} 29.0 و27.1.
\textbf{17.} $\nu_3 = \qty{38500}{cm^{-1}}$ (\qty{260}{nm}) للياقوت، في فوق البنفسجي، حيث يحجبها
امتصاص نقل الشحنة.
\textbf{18.} الموقع ثماني الأوجه في البيريل أكبر: فالروابط Cr--O الأطول تعطي انشطارًا
أصغر، يتغير بشدة مع المسافة.
\textbf{19.} $10^7/695 = \qty{14390}{cm^{-1}}$، أي $23.4B$.
\textbf{20.} يتوقف \babelsublr{$^2$E$_g$} أساسًا على $C$: فالمستوى المقيس الأعلى يعني أن $C/B$ أكبر،
نحو 5.3 بالتقطير نفسه.
\textbf{21.} ينتمي \babelsublr{$^2$E$_g$} و\babelsublr{$^4$A$_{2g}$} كلاهما إلى $t_{2g}^3$، ولا يختلفان إلا بقلب سبين: فالترابط،
ومنه الهندسة، هو نفسه في كليهما، وليس للانتقال متوالية
اهتزازية (خط 0--0).
\textbf{22.} لأنه ممنوع سبينيًا.
\textbf{23.} تتوقف طاقته على $B$ و$C$، ولا تكاد تتوقف على $\Delta_{\mathrm o}$ (خط أفقي على المخطط)،
و$B$ تكاد تكون نفسها في الحجرين.
\textbf{24.} تسترخي الأيونات المضخوخة إلى الحزم العريضة وتتراكم في المستوى الطويل العمر
\babelsublr{$^2$E$_g$}، بحيث يمكن أن يكون فيه من الأيونات أكثر مما في المستوى الأساسي: وهذا
انقلاب \omterm{def:b3:partition-functions:boltzmann}{الإشغال}.
\textbf{25.} \textbf{$\Delta_{\mathrm o}(\text{الياقوت}) - \Delta_{\mathrm o}(\text{الزمرد}) = \qty{1080}{cm^{-1}}$}: وهو كل الفرق بين الأحمر 
والأخضر.
\end{solution}
